% Options for packages loaded elsewhere
\PassOptionsToPackage{unicode}{hyperref}
\PassOptionsToPackage{hyphens}{url}
%
\documentclass[
  12pt,
  a4paper,
]{article}
\usepackage{amsmath,amssymb}
\usepackage{iftex}
\ifPDFTeX
  \usepackage[T1]{fontenc}
  \usepackage[utf8]{inputenc}
  \usepackage{textcomp} % provide euro and other symbols
\else % if luatex or xetex
  \usepackage{unicode-math} % this also loads fontspec
  \defaultfontfeatures{Scale=MatchLowercase}
  \defaultfontfeatures[\rmfamily]{Ligatures=TeX,Scale=1}
\fi
\usepackage{lmodern}
\ifPDFTeX\else
  % xetex/luatex font selection
\fi
% Use upquote if available, for straight quotes in verbatim environments
\IfFileExists{upquote.sty}{\usepackage{upquote}}{}
\IfFileExists{microtype.sty}{% use microtype if available
  \usepackage[]{microtype}
  \UseMicrotypeSet[protrusion]{basicmath} % disable protrusion for tt fonts
}{}
\makeatletter
\@ifundefined{KOMAClassName}{% if non-KOMA class
  \IfFileExists{parskip.sty}{%
    \usepackage{parskip}
  }{% else
    \setlength{\parindent}{0pt}
    \setlength{\parskip}{6pt plus 2pt minus 1pt}}
}{% if KOMA class
  \KOMAoptions{parskip=half}}
\makeatother
\usepackage{xcolor}
\usepackage[margin=1in]{geometry}
\usepackage{color}
\usepackage{fancyvrb}
\newcommand{\VerbBar}{|}
\newcommand{\VERB}{\Verb[commandchars=\\\{\}]}
\DefineVerbatimEnvironment{Highlighting}{Verbatim}{commandchars=\\\{\}}
% Add ',fontsize=\small' for more characters per line
\newenvironment{Shaded}{}{}
\newcommand{\AlertTok}[1]{\textcolor[rgb]{1.00,0.00,0.00}{\textbf{#1}}}
\newcommand{\AnnotationTok}[1]{\textcolor[rgb]{0.38,0.63,0.69}{\textbf{\textit{#1}}}}
\newcommand{\AttributeTok}[1]{\textcolor[rgb]{0.49,0.56,0.16}{#1}}
\newcommand{\BaseNTok}[1]{\textcolor[rgb]{0.25,0.63,0.44}{#1}}
\newcommand{\BuiltInTok}[1]{\textcolor[rgb]{0.00,0.50,0.00}{#1}}
\newcommand{\CharTok}[1]{\textcolor[rgb]{0.25,0.44,0.63}{#1}}
\newcommand{\CommentTok}[1]{\textcolor[rgb]{0.38,0.63,0.69}{\textit{#1}}}
\newcommand{\CommentVarTok}[1]{\textcolor[rgb]{0.38,0.63,0.69}{\textbf{\textit{#1}}}}
\newcommand{\ConstantTok}[1]{\textcolor[rgb]{0.53,0.00,0.00}{#1}}
\newcommand{\ControlFlowTok}[1]{\textcolor[rgb]{0.00,0.44,0.13}{\textbf{#1}}}
\newcommand{\DataTypeTok}[1]{\textcolor[rgb]{0.56,0.13,0.00}{#1}}
\newcommand{\DecValTok}[1]{\textcolor[rgb]{0.25,0.63,0.44}{#1}}
\newcommand{\DocumentationTok}[1]{\textcolor[rgb]{0.73,0.13,0.13}{\textit{#1}}}
\newcommand{\ErrorTok}[1]{\textcolor[rgb]{1.00,0.00,0.00}{\textbf{#1}}}
\newcommand{\ExtensionTok}[1]{#1}
\newcommand{\FloatTok}[1]{\textcolor[rgb]{0.25,0.63,0.44}{#1}}
\newcommand{\FunctionTok}[1]{\textcolor[rgb]{0.02,0.16,0.49}{#1}}
\newcommand{\ImportTok}[1]{\textcolor[rgb]{0.00,0.50,0.00}{\textbf{#1}}}
\newcommand{\InformationTok}[1]{\textcolor[rgb]{0.38,0.63,0.69}{\textbf{\textit{#1}}}}
\newcommand{\KeywordTok}[1]{\textcolor[rgb]{0.00,0.44,0.13}{\textbf{#1}}}
\newcommand{\NormalTok}[1]{#1}
\newcommand{\OperatorTok}[1]{\textcolor[rgb]{0.40,0.40,0.40}{#1}}
\newcommand{\OtherTok}[1]{\textcolor[rgb]{0.00,0.44,0.13}{#1}}
\newcommand{\PreprocessorTok}[1]{\textcolor[rgb]{0.74,0.48,0.00}{#1}}
\newcommand{\RegionMarkerTok}[1]{#1}
\newcommand{\SpecialCharTok}[1]{\textcolor[rgb]{0.25,0.44,0.63}{#1}}
\newcommand{\SpecialStringTok}[1]{\textcolor[rgb]{0.73,0.40,0.53}{#1}}
\newcommand{\StringTok}[1]{\textcolor[rgb]{0.25,0.44,0.63}{#1}}
\newcommand{\VariableTok}[1]{\textcolor[rgb]{0.10,0.09,0.49}{#1}}
\newcommand{\VerbatimStringTok}[1]{\textcolor[rgb]{0.25,0.44,0.63}{#1}}
\newcommand{\WarningTok}[1]{\textcolor[rgb]{0.38,0.63,0.69}{\textbf{\textit{#1}}}}
\setlength{\emergencystretch}{3em} % prevent overfull lines
\providecommand{\tightlist}{%
  \setlength{\itemsep}{0pt}\setlength{\parskip}{0pt}}
\setcounter{secnumdepth}{-\maxdimen} % remove section numbering
% definitions for citeproc citations
\NewDocumentCommand\citeproctext{}{}
\NewDocumentCommand\citeproc{mm}{%
  \begingroup\def\citeproctext{#2}\cite{#1}\endgroup}
\makeatletter
 % allow citations to break across lines
 \let\@cite@ofmt\@firstofone
 % avoid brackets around text for \cite:
 \def\@biblabel#1{}
 \def\@cite#1#2{{#1\if@tempswa , #2\fi}}
\makeatother
\newlength{\cslhangindent}
\setlength{\cslhangindent}{1.5em}
\newlength{\csllabelwidth}
\setlength{\csllabelwidth}{3em}
\newenvironment{CSLReferences}[2] % #1 hanging-indent, #2 entry-spacing
 {\begin{list}{}{%
  \setlength{\itemindent}{0pt}
  \setlength{\leftmargin}{0pt}
  \setlength{\parsep}{0pt}
  % turn on hanging indent if param 1 is 1
  \ifodd #1
   \setlength{\leftmargin}{\cslhangindent}
   \setlength{\itemindent}{-1\cslhangindent}
  \fi
  % set entry spacing
  \setlength{\itemsep}{#2\baselineskip}}}
 {\end{list}}
\usepackage{calc}
\newcommand{\CSLBlock}[1]{\hfill\break\parbox[t]{\linewidth}{\strut\ignorespaces#1\strut}}
\newcommand{\CSLLeftMargin}[1]{\parbox[t]{\csllabelwidth}{\strut#1\strut}}
\newcommand{\CSLRightInline}[1]{\parbox[t]{\linewidth - \csllabelwidth}{\strut#1\strut}}
\newcommand{\CSLIndent}[1]{\hspace{\cslhangindent}#1}
\usepackage{fvextra}\DefineVerbatimEnvironment{verbatim}{Verbatim}{breaklines}\DefineVerbatimEnvironment{Highlighting}{Verbatim}{commandchars=\\\{\},breaklines}
\ifLuaTeX
  \usepackage{selnolig}  % disable illegal ligatures
\fi
\usepackage{bookmark}
\IfFileExists{xurl.sty}{\usepackage{xurl}}{} % add URL line breaks if available
\urlstyle{same}
\hypersetup{
  hidelinks,
  pdfcreator={LaTeX via pandoc}}

\author{}
\date{}

\begin{document}

\section{Build a Digital Twin for a Potted
Plant}\label{build-a-digital-twin-for-a-potted-plant}

\textbf{What you'll build.} A \textasciitilde70-line Python script,
\texttt{twin.py}, that simulates a potted plant losing moisture over 24
hours and a digital twin that reads its moisture sensor, predicts when
the soil will dry out, and waters the plant automatically -- no human in
the loop. Running it prints a line per hour and, partway through, one
marked \texttt{WATERED}:

\begin{verbatim}
hour 11  sensor= 38.2  predicted_next= 35.0
hour 12  sensor= 36.3  predicted_next= 34.4  WATERED
hour 13  sensor= 59.3  predicted_next= 82.3
\end{verbatim}

\textbf{What you'll learn.} How to structure a digital twin as two
independent pieces -- a physical-side simulation and a twin that only
ever sees sensor readings -- and why closing the loop (the twin issuing
a command back to the physical side) is what separates a digital twin
from a dashboard.

\textbf{What you need.} Python 3.10 or later, and a terminal. No
packages beyond the standard library.

\textbf{Time.} About 20 minutes.

\subsection{Step 1: Create the project
folder}\label{step-1-create-the-project-folder}

\begin{verbatim}
mkdir plant-twin
cd plant-twin
\end{verbatim}

\subsection{Step 2: Simulate the pot}\label{step-2-simulate-the-pot}

You don't have a real plant wired up, so start by writing a stand-in for
one. Create \texttt{twin.py} with this content:

\begin{Shaded}
\begin{Highlighting}[]
\ImportTok{import}\NormalTok{ random}

\NormalTok{random.seed(}\DecValTok{7}\NormalTok{)}


\KeywordTok{class}\NormalTok{ Pot:}
    \CommentTok{"""Stands in for the real pot\textquotesingle{}s soil {-}{-} in a hardware deployment this}
\CommentTok{    state lives in wet dirt, not a Python object."""}

    \KeywordTok{def} \FunctionTok{\_\_init\_\_}\NormalTok{(}\VariableTok{self}\NormalTok{, moisture}\OperatorTok{=}\FloatTok{60.0}\NormalTok{):}
        \VariableTok{self}\NormalTok{.moisture }\OperatorTok{=}\NormalTok{ moisture}

    \KeywordTok{def}\NormalTok{ tick(}\VariableTok{self}\NormalTok{):}
        \VariableTok{self}\NormalTok{.moisture }\OperatorTok{{-}=}\NormalTok{ random.uniform(}\FloatTok{1.5}\NormalTok{, }\FloatTok{2.5}\NormalTok{)}
        \VariableTok{self}\NormalTok{.moisture }\OperatorTok{=} \BuiltInTok{max}\NormalTok{(}\FloatTok{0.0}\NormalTok{, }\VariableTok{self}\NormalTok{.moisture)}

    \KeywordTok{def}\NormalTok{ water(}\VariableTok{self}\NormalTok{):}
        \VariableTok{self}\NormalTok{.moisture }\OperatorTok{=} \BuiltInTok{min}\NormalTok{(}\FloatTok{100.0}\NormalTok{, }\VariableTok{self}\NormalTok{.moisture }\OperatorTok{+} \FloatTok{25.0}\NormalTok{)}

    \KeywordTok{def}\NormalTok{ read\_sensor(}\VariableTok{self}\NormalTok{):}
        \ControlFlowTok{return} \BuiltInTok{round}\NormalTok{(}\VariableTok{self}\NormalTok{.moisture }\OperatorTok{+}\NormalTok{ random.uniform(}\OperatorTok{{-}}\FloatTok{1.0}\NormalTok{, }\FloatTok{1.0}\NormalTok{), }\DecValTok{1}\NormalTok{)}
\end{Highlighting}
\end{Shaded}

\texttt{tick()} is one hour passing: the soil loses a random amount of
moisture, the way real soil does under evaporation.
\texttt{read\_}\penalty0\texttt{sensor()} is what a real capacitive
soil-moisture probe gives you -- the true value, plus a little sensor
noise. \texttt{random.seed(7)} pins that randomness so your output
matches this tutorial's exactly.

Check that it works before moving on:

\begin{verbatim}
python3 -c "
from twin import Pot
pot = Pot()
for _ in range(3):
    pot.tick()
    print(pot.read_sensor())
"
\end{verbatim}

You should see:

\begin{verbatim}
57.5
55.2
53.7
\end{verbatim}

Three ticks, three shrinking readings. That's the physical side done --
notice it knows nothing about digital twins. It's just a pot.

\subsection{Step 3: Write the digital
twin}\label{step-3-write-the-digital-twin}

The twin is a second, separate class. Critically, it never touches
\texttt{Pot} directly -- it only ever sees numbers arriving from
\texttt{read\_}\penalty0\texttt{sensor()}, the same way a real twin only
sees whatever its sensor's network connection delivers. Append this to
\texttt{twin.py}:

\begin{Shaded}
\begin{Highlighting}[]


\KeywordTok{class}\NormalTok{ PlantTwin:}
    \CommentTok{"""The digital twin. It never touches Pot directly {-}{-} only sensor}
\CommentTok{    readings flow in, and water commands flow back out."""}

\NormalTok{    DRY\_THRESHOLD }\OperatorTok{=} \FloatTok{35.0}

    \KeywordTok{def} \FunctionTok{\_\_init\_\_}\NormalTok{(}\VariableTok{self}\NormalTok{):}
        \VariableTok{self}\NormalTok{.history }\OperatorTok{=}\NormalTok{ []}

    \KeywordTok{def}\NormalTok{ ingest(}\VariableTok{self}\NormalTok{, reading):}
        \VariableTok{self}\NormalTok{.history.append(reading)}
        \VariableTok{self}\NormalTok{.history }\OperatorTok{=} \VariableTok{self}\NormalTok{.history[}\OperatorTok{{-}}\DecValTok{3}\NormalTok{:]}

    \KeywordTok{def}\NormalTok{ predict\_next(}\VariableTok{self}\NormalTok{):}
        \ControlFlowTok{if} \BuiltInTok{len}\NormalTok{(}\VariableTok{self}\NormalTok{.history) }\OperatorTok{\textless{}} \DecValTok{2}\NormalTok{:}
            \ControlFlowTok{return} \VariableTok{self}\NormalTok{.history[}\OperatorTok{{-}}\DecValTok{1}\NormalTok{]}
\NormalTok{        trend }\OperatorTok{=} \VariableTok{self}\NormalTok{.history[}\OperatorTok{{-}}\DecValTok{1}\NormalTok{] }\OperatorTok{{-}} \VariableTok{self}\NormalTok{.history[}\OperatorTok{{-}}\DecValTok{2}\NormalTok{]}
        \ControlFlowTok{return} \VariableTok{self}\NormalTok{.history[}\OperatorTok{{-}}\DecValTok{1}\NormalTok{] }\OperatorTok{+}\NormalTok{ trend}

    \KeywordTok{def}\NormalTok{ decide(}\VariableTok{self}\NormalTok{):}
        \ControlFlowTok{return} \VariableTok{self}\NormalTok{.predict\_next() }\OperatorTok{\textless{}} \VariableTok{self}\NormalTok{.DRY\_THRESHOLD}
\end{Highlighting}
\end{Shaded}

\texttt{ingest()} is the twin's half of the data connection: each new
reading gets appended, and only the last three are kept, since only the
recent trend matters for what comes next.
\texttt{predict\_}\penalty0\texttt{next()} is the twin's model --
deliberately the simplest one that works: extend the line between the
last two readings one more hour forward. \texttt{decide()} is the twin's
algorithm: if that predicted value would be too dry, it's time to water.
This is exactly the three-part shape every digital twin shares -- data
in, a model to interpret it, an algorithm to act on it -- just small
enough here to read in one sitting.

\subsection{Step 4: Close the loop}\label{step-4-close-the-loop}

A twin that only predicts is a dashboard. What makes it a \emph{twin} is
that its decision changes the physical side. Append the loop that wires
\texttt{Pot} and \texttt{PlantTwin} together:

\begin{Shaded}
\begin{Highlighting}[]


\KeywordTok{def}\NormalTok{ run(hours):}
\NormalTok{    pot }\OperatorTok{=}\NormalTok{ Pot()}
\NormalTok{    twin }\OperatorTok{=}\NormalTok{ PlantTwin()}

    \ControlFlowTok{for}\NormalTok{ hour }\KeywordTok{in} \BuiltInTok{range}\NormalTok{(}\DecValTok{1}\NormalTok{, hours }\OperatorTok{+} \DecValTok{1}\NormalTok{):}
\NormalTok{        pot.tick()}
\NormalTok{        reading }\OperatorTok{=}\NormalTok{ pot.read\_sensor()}
\NormalTok{        twin.ingest(reading)}

\NormalTok{        watered }\OperatorTok{=} \VariableTok{False}
        \ControlFlowTok{if}\NormalTok{ twin.decide():}
\NormalTok{            pot.water()}
\NormalTok{            watered }\OperatorTok{=} \VariableTok{True}

\NormalTok{        status }\OperatorTok{=} \StringTok{"WATERED"} \ControlFlowTok{if}\NormalTok{ watered }\ControlFlowTok{else} \StringTok{""}
        \BuiltInTok{print}\NormalTok{(}\SpecialStringTok{f"hour }\SpecialCharTok{\{}\NormalTok{hour}\SpecialCharTok{:2d\}}\SpecialStringTok{  sensor=}\SpecialCharTok{\{}\NormalTok{reading}\SpecialCharTok{:5.1f\}}\SpecialStringTok{  "}
              \SpecialStringTok{f"predicted\_next=}\SpecialCharTok{\{}\NormalTok{twin}\SpecialCharTok{.}\NormalTok{predict\_next()}\SpecialCharTok{:5.1f\}}\SpecialStringTok{  }\SpecialCharTok{\{}\NormalTok{status}\SpecialCharTok{\}}\SpecialStringTok{"}\NormalTok{)}


\ControlFlowTok{if} \VariableTok{\_\_name\_\_} \OperatorTok{==} \StringTok{"\_\_main\_\_"}\NormalTok{:}
\NormalTok{    run(}\DecValTok{24}\NormalTok{)}
\end{Highlighting}
\end{Shaded}

Each hour, the pot's real moisture drops, the sensor reports a noisy
version of it, the twin folds that reading into its prediction, and if
the twin's prediction crosses the dry threshold, it calls
\texttt{pot.water()} -- reaching back across the loop to change the
physical side, before a human ever looks at the readings.

\subsection{Step 5: Run it}\label{step-5-run-it}

\begin{verbatim}
python3 twin.py
\end{verbatim}

You should see 24 lines of output. Line 12 is marked \texttt{WATERED},
and the next line's \texttt{sensor} value jumps back up above 59,
because the twin's command already changed the pot before that reading
was taken:

\begin{verbatim}
hour 11  sensor= 38.2  predicted_next= 35.0
hour 12  sensor= 36.3  predicted_next= 34.4  WATERED
hour 13  sensor= 59.3  predicted_next= 82.3
\end{verbatim}

If your output matches, the loop is genuinely closed: the twin's
decision at hour 12 is the reason hour 13's real-world reading looks
completely different.

\subsection{Step 6: Change the twin's behavior without touching the
pot}\label{step-6-change-the-twins-behavior-without-touching-the-pot}

Open \texttt{twin.py} and change:

\begin{Shaded}
\begin{Highlighting}[]
\NormalTok{    DRY\_THRESHOLD }\OperatorTok{=} \FloatTok{35.0}
\end{Highlighting}
\end{Shaded}

to:

\begin{Shaded}
\begin{Highlighting}[]
\NormalTok{    DRY\_THRESHOLD }\OperatorTok{=} \FloatTok{45.0}
\end{Highlighting}
\end{Shaded}

Save the file and run \texttt{python3}\penalty0\texttt{\ twin.py} again.
Watering now happens at hour 8 instead of hour 12, and a second time at
hour 21:

\begin{verbatim}
hour  8  sensor= 45.2  predicted_next= 42.3  WATERED
...
hour 21  sensor= 44.5  predicted_next= 42.0  WATERED
\end{verbatim}

Notice what you didn't have to change: \texttt{Pot} is untouched. This
is the payoff of keeping the twin as a separate piece that only talks to
the physical side through readings and commands -- you can retune,
replace, or completely rewrite the twin's model without ever touching
the thing it's twinning.

\subsection{What you built}\label{what-you-built}

You now have a working digital twin: a physical-side simulation that
knows nothing about twins, a twin that knows nothing about the physical
side's internals, and a closed loop connecting them where the twin's own
prediction changes what the physical side does next. Swap \texttt{Pot}
for code that reads a real capacitive soil sensor over I2C and calls a
relay to run a water pump, and the \texttt{PlantTwin} class doesn't need
to change at all -- that's the whole reason to build the two halves this
way.

\subsection{Where to go next}\label{where-to-go-next}

\texttt{PlantTwin}'s ``extend the last two points'' model is
deliberately the crudest thing that could work. A real deployment would
replace \texttt{predict\_}\penalty0\texttt{next} with a Kalman filter,
which blends the model's prediction against each new noisy reading
instead of trusting the latest one alone {[}1{]}, and would read real
sensor data over MQTT rather than calling a Python method directly. A
how-to guide would be the right place to cover wiring up an actual
moisture probe and pump, since that involves choices -- which board,
which protocol -- that a first tutorial deliberately avoids.

\subsection{References}\label{references}

\phantomsection\label{refs}
\begin{CSLReferences}{0}{0}
\bibitem[\citeproctext]{ref-feng_model-based_2023}
\CSLLeftMargin{{[}1{]} }%
\CSLRightInline{H. Feng, C. Gomes, and P. G. Larsen, {``Model-{Based}
{Monitoring} and {State} {Estimation} for {Digital} {Twins}: {The}
{Kalman} {Filter}.''} arXiv, Apr. 2023. doi:
\href{https://doi.org/10.48550/arXiv.2305.00252}{10.48550/arXiv.2305.00252}.}

\end{CSLReferences}

\end{document}
